% 20260703_Assingment_Probability_Statistics_A
\documentclass[dvipdfmx]{jsreport}
%preamble
\usepackage{soupreamble}

\title{確率及び統計学特論A 第5回レポート課題}
\author{M1 6012 川村聡一郎}
\date{}


\begin{document}
\maketitle

\begin{tcolorbox}
	$\mu$の分布関数を$F$とし、$F$の不連続点を$A_\mu$とする。このとき$A_\mu$は高々加算であることを示せ。
%Hint: 分布関数をFとして、$A_{\mu, n}= \{x\in \R; F(x)- F(x-)> 1/n\}$とおくと
% #(A_{\mu, n}の元)\le n(背理法)
\end{tcolorbox}

\begin{souanswer} %答案はここに書く
\begin{enumerate}
	\item 
\end{enumerate}
\end{souanswer}
\end{document}